\section{遗址、碑石与后人的金：材料怎样留下北方}

\subsection{一块碑石不能代替一座城}

金代碑刻常给人一种接近地方生活的感觉：它有姓名、官衔、日期、籍贯和捐施事项，似乎比帝纪更具体。具体正是它的价值，但也是它的界限。立碑者选择要被记住的祖先、功德与关系，石料能够保存下来又取决于战乱、迁徙、再利用和近代收藏。一通碑最可靠地说明的是它所记的人和事在特定地点进入了纪念；它不能代表同城未立碑者，更不能变成全朝社会的抽样。

墓志也如此。它们常保存迁居、仕宦、婚姻与家族叙述，能使抽象的“人口流动”落到一个家庭的时间上；同时，墓志以精英的价值判断组织生平，追赠官和美化的德行须同实授、任所和成文场景分开。将墓志与正史、地方地理和物质遗存对读，可以辨认某些迁徙或任官的轮廓，却不应把纪念文字变成透明的私人日记。

\subsection{中都遗址告诉我们的尺度}

中都的城垣、宫城、道路和作坊遗迹，使金朝都城不再只是《金史》中一个迁都年份。它们让研究者追问建筑怎样布置、城市空间怎样被使用、工程需要何种材料。遗址也有明确的发掘边界：已揭露的区域不等于整座城市，后代扰动会改变层位，器物的生产、使用和埋藏时间未必一致。考古提供的是物质活动的证据，不是能自动填满人物动机的舞台。

将都城遗址放回文本中，能够改变对行政的理解。迁都不只是君主把宫廷从一地搬到另一地，而是要重组水陆运输、官署、市场、劳动力和防卫。文本能标出命令和工程，遗址能限制空间想象；二者相遇之处，才可以讨论金廷怎样把华北资源集中到中都。二者没有说出的地方，则不能用华丽的都城场景补造。

\subsection{器物的旅行与人的旅行}

钱币、陶瓷、金属器和经卷会离开生产地，进入市场、墓葬、窖藏或寺院。它们的移动说明区域之间存在交换或接触，却不必等于一条由国家控制的贸易线。窖藏可能因战乱被埋，也可能有其他原因；一件外来器物可能由商旅带来，也可能经过多次转手。材料的旅行路线常比人的旅行更容易保存，但不能让器物代替那些搬运、购买和使用它的人。

这类遗存仍然重要，因为它们使北方不再只按王朝边界理解。东北、燕云、山东、河南和边区通过货物、纸张、工艺和宗教物品相连；战争会切断一些线路，也会让军需创造新的流动。地方经济不是帝纪的背景，而是政权能否获得粮、税和信息的条件。考古与文献相互校正，能让这种条件从抽象的“经济基础”变为实际可追问的道路、作坊和仓场。

\subsection{谁编写了金史}

《金史》成书于元末，保存了金朝国史与档案传统的大量片段，也以元朝修史的秩序整理前朝。它对于帝纪、制度名目、官员履历和战事年序不可或缺；其叙事重点仍在皇帝、中枢和可进入国史的人物。宋方的编年与笔记从敌对或相邻政权的立场观察金，常保存使节、边界和战争细节，却有自身的政治压力。两类文字并列，不会自然产生中立答案，却能使单一王朝的盲点显露。

后来的地方志、家谱和文学作品还要更谨慎。它们可能保留早期地名、家族记忆或文化想象，也可能把晚出的制度投回金代，把后人的需要写成祖先的事实。使用它们的正确方式不是一概排除，而是辨别成书时间、文本目的和能否同较近材料互证。史料批判不是在叙事外加一段免责声明；它决定哪些细节能被写成确定，哪些只能作为可能。

\subsection{金亡之后的记忆}

金亡并未使女真、契丹、汉人或其他北方居民的记忆同时消失。后来的政权重新安排官名、道路和税役，家族则可能以迁居、仕宦或宗教捐施叙述自身过去。元、明以后的文字会以不同政治问题回望金：有人强调征服，有人强调北方秩序，有人以宋金对峙组织道德判断。读者在这些回望中见到的，既是金代留下的痕迹，也是后人需要一个怎样的金。

因此，本章的终点不在于给金朝贴上成功或失败的总标签。东北起兵、接收燕云、经营华北、迁都中都、维持南界和抵抗蒙古，都让不同人进入新的道路、税役、文书与记忆关系。王朝的终结使这些关系被重新分配，却不能使它们从历史中撤去。把正史、宋方材料、碑志和考古并置，正是为了让这种不整齐的延续保持可见。
